Grundsätzlich sollte einfacher Zeilenabstand zu Anwendung kommen. Um Platz
für Anmekrungen des Gutachters zu bieten, kann der Zeilenabstand auch auf
1\nicefrac{1}{2}-fachen Zeilenabstand vergrößert werden. Bei noch größerem
Zeilenabstand könnte der Eindruck entstehen, Sie wollen die Arbeit
umfangreicher erscheinen lassen, als sie ist.

\todo[inline]{Code-Beispiel, Paket \texttt{setspace}. Normaler Zeilenabstand in Fußnoten, Minipages, Gleitumgebungen}



%\setstretch{1.25}

%% setspace benutzt in Gleitumgebungen den normalen Zeilenabstand. Das wollen wir nicht.
%\makeatletter
%\let\@xfloat=\latex@xfloat
%\makeatother
%
%% Auch in den Fußnoten beim größeren Zeilenabstand bleiben.
%\makeatletter
%\long\def\@footnotetext#1{%
%  \insert\footins{%
%    \reset@font\footnotesize
%    \interlinepenalty\interfootnotelinepenalty
%    \splittopskip\footnotesep
%    \splitmaxdepth \dp\strutbox \floatingpenalty \@MM
%    \hsize\columnwidth
%    \@parboxrestore
%    \protected@edef\@currentlabel{%
%      \csname p@footnote\endcsname\@thefnmark
%    }%    
%    \color@begingroup
%      \@makefntext{%
%        \rule\z@\footnotesep\ignorespaces#1\@finalstrut\strutbox}%
%    \color@endgroup}}
%\makeatother
%
%\makeatletter
%\long\def\@mpfootnotetext#1{%
%  \global\setbox\@mpfootins\vbox{%
%    \unvbox \@mpfootins
%    \reset@font\footnotesize
%    \hsize\columnwidth
%    \@parboxrestore
%    \protected@edef\@currentlabel{%
%      \csname p@mpfootnote\endcsname\@thefnmark}%
%    \color@begingroup
%      \@makefntext{%
%       \rule\z@\footnotesep\ignorespaces#1\@finalstrut\strutbox}%
%   \color@endgroup}}
%\makeatother
